\documentclass[11pt]{article}
\usepackage[T1]{fontenc}
\usepackage[utf8]{inputenc}
\usepackage{amsmath,amssymb,amsthm,geometry,booktabs}
\usepackage[hidelinks]{hyperref}
\geometry{margin=1in}
\newtheorem{proposition}{Proposition}
\newtheorem{remark}{Remark}
\title{Ramanujan, Landen, and Baker N-Series Refinement:\\Small-Nome Gap Asymptotics and a Two-Axis Comparison}
\author{Wayne Baker}
\date{Original: June 2026\\Review revision: October 24, 2026}
\begin{document}\maketitle
\begin{abstract}
We compare three distinct refinement mechanisms: Ramanujan's ellipse-perimeter approximation along the eccentricity axis, fixed-scale geometric N-series cancellation along the resolution axis, and Landen--arithmetic--geometric mean (AGM) contraction along the modulus axis. The classical theta-function identity $\delta(q)=k(q^2)$, for the normalized complementary-modulus gap $\delta(q)=(1-k'(q))/(1+k'(q))$, directly yields, for every fixed integer $m\ge2$,
\[
\frac{\delta(q^m)}{\delta(q)^m}=4^{1-m}\left(1+4m q^2+O(q^4)\right),\qquad q\to0.
\]
The first-correction coefficient is thus linear in $m$; this is a small-nome asymptotic, not a newly derived degree-$m$ modular equation. With explicit Gauss--Legendre indexing, classical AGM error asymptotics imply $\alpha_n/\delta_n^2\to1/8$ and $\alpha_{n+1}/\alpha_n^2\to1/2$. We separate these specific coefficients from generic normalized quadratic and cubic error maps. We make no claim that geometric N-series extrapolation is a modular transformation.
\end{abstract}

\section{Ramanujan's approximation: the eccentricity axis}
For ellipse semiaxes $a\ge b>0$, set $h=((a-b)/(a+b))^2$. The exact perimeter is
\[
P=\pi(a+b)\,{}_2F_1(-\tfrac12,-\tfrac12;1;h)
=\pi(a+b)\left(1+\tfrac14h+\tfrac1{64}h^2+\tfrac1{256}h^3+\tfrac{25}{16384}h^4+\cdots\right).
\]
Ramanujan's second approximation is
\[
R_2=\pi(a+b)\left(1+\frac{3h}{10+\sqrt{4-3h}}\right).
\]
It agrees through $h^4$, and our \emph{signed-error convention} is
\[
\boxed{P-R_2=\pi(a+b)\left(\frac{3}{131072}h^5+O(h^6)\right).}
\]
Equivalently, the coefficient in $R_2-P$ is $-3/131072$. This mechanism is an eccentricity expansion; it does not refine a polygonal resolution variable.

\section{Geometric N-series: the resolution axis}
If a polygonal perimeter family admits an even-power expansion of suitable order,
\[
P_N=P+C_2N^{-2}+C_4N^{-4}+C_6N^{-6}+\cdots,
\]
Richardson-type linear combinations of $P_N,P_{rN},\ldots$ may cancel selected terms. Cancelling the first $s$ error terms yields $O(N^{-2(s+1)})$ provided the necessary expansion and remainder bound hold. The cancellation coefficients arise from a finite Vandermonde/moment system. This algebraic convergence in $N$ is distinct from nome-squaring under the descending Landen transformation and from the cubic trisection residual. No fixed number of Richardson columns creates a modular transformation.

\section{Landen contraction and explicit AGM indexing}
Define the Gauss--Legendre iteration by
\begin{align*}
a_0&=1,&b_0&=2^{-1/2},&t_0&=\tfrac14,&p_0&=1,\\
a_{n+1}&=\tfrac12(a_n+b_n),&b_{n+1}&=\sqrt{a_nb_n},
&t_{n+1}&=t_n-p_n(a_n-a_{n+1})^2,&p_{n+1}&=2p_n.
\end{align*}
We use \emph{matching indices} throughout:
\[
\boxed{\pi_n=\frac{(a_n+b_n)^2}{4t_n},\quad
\delta_n=\frac{a_n-b_n}{a_n+b_n},\quad e_n=\pi_n-\pi,\quad
\alpha_n=-\frac{e_{n+1}}{e_{n+1}-e_n}.}
\]
In particular $\pi_n$ is formed from $(a_n,b_n,t_n)$, not from the next iteration's data. Altering this convention alters intermediate identities and may alter reported constants. The AGM update gives the exact descending Landen relation
\[
\delta_{n+1}=\frac{1-\sqrt{1-\delta_n^2}}{1+\sqrt{1-\delta_n^2}}
=\frac{\delta_n^2}{(1+\sqrt{1-\delta_n^2})^2}
=\frac14\delta_n^2+O(\delta_n^4).
\]
For the above Gauss--Legendre convention, the classical error asymptotics can be expressed as
\[
e_n\sim-\pi^2 2^n\delta_n^2.
\]
Consequently $|e_{n+1}|/|e_n|\to0$, and the exact definition of $\alpha_n$ gives
\[
\alpha_n\sim\frac{e_{n+1}}{e_n}
\sim2\frac{\delta_{n+1}^2}{\delta_n^2}
\sim\frac18\delta_n^2,
\qquad\boxed{\lim_{n\to\infty}\frac{\alpha_n}{\delta_n^2}=\frac18.}
\]
The same relations imply
\[
\boxed{\lim_{n\to\infty}\frac{\alpha_{n+1}}{\alpha_n^2}=\frac12.}
\]
These limiting constants use the stated definitions; they are not invariant under arbitrary shifts in the approximate-$\pi$ indexing.

\section{Small-nome gap asymptotics, not higher-degree modular equations}
Let $q=e^{i\pi\tau}$, $\operatorname{Im}\tau>0$, and take $0<q<1$ in the real-valued statements. Write
\[
k(q)=\frac{\vartheta_2(q)^2}{\vartheta_3(q)^2},\quad
k'(q)=\sqrt{1-k(q)^2},\quad
\delta(q)=\frac{1-k'(q)}{1+k'(q)}.
\]
The classical quadratic theta identities give the exact identity
\begin{equation}\label{eq:delta_nome}
\boxed{\delta(q)=k(q^2).}
\end{equation}
Indeed $\vartheta_3(q)^2+\vartheta_4(q)^2=2\vartheta_3(q^2)^2$ and
$\vartheta_3(q)^2-\vartheta_4(q)^2=2\vartheta_2(q^2)^2$; since $k'(q)=\vartheta_4(q)^2/\vartheta_3(q)^2$, their quotient yields \eqref{eq:delta_nome}.
The classical theta product then yields
\[
k(Q)=4\sqrt{Q}\prod_{j=1}^\infty\left(\frac{1+Q^{2j}}{1+Q^{2j-1}}\right)^4,
\quad\delta(q)=k(q^2)=4q\left(1-4q^2+O(q^4)\right).
\]
\begin{proposition}[Fixed-degree small-nome gap ratio]\label{prop:gap}
For every \emph{fixed} integer $m\ge2$, as $q\to0^+$,
\[
\boxed{\frac{\delta(q^m)}{\delta(q)^m}
=4^{1-m}\left(1+4m q^2+O(q^4)\right).}
\]
In particular, the ratio tends to $4^{1-m}$.
\end{proposition}
\begin{proof}
By \eqref{eq:delta_nome}, $\delta(q^m)=k(q^{2m})=4q^m(1-4q^{2m}+O(q^{4m}))$ and $\delta(q)^m=4^mq^m(1-4q^2+O(q^4))^m$. Divide, expand the inverse denominator, and note that $2m\ge4$. The coefficient of $q^2$ is therefore $4m$.
\end{proof}
\begin{remark}[Terminology and priority]
This statement concerns the leading and first-correction asymptotics of $k(q^{2m})$ relative to $k(q^2)^m$. It is \emph{not} the classical degree-$m$ modular polynomial or a new modular equation. The ingredients are standard theta-function identities and products; no historical priority is claimed for them.
\end{remark}

\section{Numerical illustration}
For $R_m(q)=\delta(q^m)/\delta(q)^m$, define $C_m(q)=q^{-2}(R_m(q)/4^{1-m}-1)$. The proposition predicts $C_m(q)\to4m$. At $q=10^{-12}$ with 400-digit working arithmetic, the reported values are:
\begin{center}\begin{tabular}{rrl}\toprule
$m$&$4m$&$C_m(10^{-12})$\\\midrule
2&8&8.000000000000000000000016\\
3&12&12.000000000000000000000054\\
4&16&16.000000000000000000000104\\
5&20&20.000000000000000000000170\\
6&24&24.000000000000000000000252\\\bottomrule
\end{tabular}\end{center}
These are numerical illustrations, not the proof. To avoid cancellation when $q$ is tiny, use $\delta(q)=\lambda(q)/(1+\sqrt{1-\lambda(q)})^2$, where $\lambda(q)=k(q)^2$. A canonical repository command is
\begin{center}
\begin{minipage}{0.92\linewidth}
\begin{verbatim}
python scripts/verify_modular_gap_law.py --precision 400 --outdir outputs
\end{verbatim}
\end{minipage}
\end{center}

\section{What normalization does and does not show}
For a local error law $x_{n+1}=\lambda x_n^p+O(x_n^{p+2})$ with $\lambda\ne0$, constant rescaling can normalize the leading coefficient to $\pm1$; hence the bare normalized form primarily encodes the local \emph{order} $p$. In particular, the Landen law is $\delta_{n+1}=\tfrac14\delta_n^2+O(\delta_n^4)$, while with $\eta_n=\delta_n/4$ it becomes $\eta_{n+1}=\eta_n^2+O(\eta_n^4)$. These equations encode the same quadratic contraction.

The companion Baker $4{:}3$ trisection analysis records a cubic residual and, in its specified geometric normalization, the coefficient $-7/384$. Where its signed error variable $x_n$ obeys $x_{n+1}=-(7/384)x_n^3+O(x_n^5)$, putting $\eta_n=\sqrt{7/384}\,x_n$ gives $\eta_{n+1}=-\eta_n^3+O(\eta_n^5)$. The value $7/384$ belongs to the \emph{chosen geometric residual}, whereas the normalized equation by itself cannot recover that value. The comparison between degree two and degree three is therefore a comparison of convergence orders, not a claim of a shared special coefficient or of modular equivalence.

\section*{AI-assistance and review disclosure}
ChatGPT and Claude were used for drafting assistance, mathematical checking,
numerical diagnostics, and editorial review. High-precision agreement is
computational evidence, not independent human peer review. The analytical
claims remain subject to the stated hypotheses and cited classical results.
\section{Conclusion and limits}
Ramanujan's second ellipse approximation achieves unusually high-order cancellation in the eccentricity parameter $h$; the geometric N-series removes specified even powers in polygonal resolution $N^{-1}$; and the AGM/Landen iteration contracts a modulus-like gap quadratically. The exact theta identity $\delta(q)=k(q^2)$ streamlines the small-nome ratio proof, clarifies the nomenclature, and situates the result within classical modular theory. The one-eighth AGM constant is tied to explicitly indexed iterates, and the cubic trisection coefficient cannot be inferred from a normal form alone. These mechanisms are comparable as forms of error suppression, but are not mathematically identical.

\begin{thebibliography}{99}
\bibitem{berndt} B.~C. Berndt, \emph{Ramanujan's Notebooks, Part III}, Springer, 1991.
\bibitem{borwein} J.~M. Borwein and P.~B. Borwein, \emph{Pi and the AGM}, Wiley, 1987.
\bibitem{brent} R.~P. Brent, ``Fast multiple-precision evaluation of elementary functions,'' \emph{J. ACM} 23 (1976), 242--251.
\bibitem{watson} E.~T. Whittaker and G.~N. Watson, \emph{A Course of Modern Analysis}, 4th ed., Cambridge, 1927.
\bibitem{lawden} D.~F. Lawden, \emph{Elliptic Functions and Applications}, Springer, 1989.
\bibitem{abramowitz} M.~Abramowitz and I.~A. Stegun (eds.), \emph{Handbook of Mathematical Functions}, Dover, 1972.
\bibitem{richardson} L.~F. Richardson, ``The approximate arithmetical solution by finite differences of physical problems,'' \emph{Phil. Trans. R. Soc. A} 210 (1911), 307--357.
\bibitem{trefethen} L.~N. Trefethen and J.~A.~C. Weideman, ``The exponentially convergent trapezoidal rule,'' \emph{SIAM Review} 56 (2014), 385--458.
\bibitem{bakeragm} W.~Baker, \emph{A Prefactor Coefficient Law for the Gauss--Legendre AGM Algorithm}, 2026, \url{https://doi.org/10.5281/zenodo.22896047}.
\bibitem{bakern} W.~Baker, \emph{A Constructive N-Series Acceleration Law for Polygonal Approximation of $\pi$}, preprint, 2026.
\bibitem{bakerc} W.~Baker, \emph{A Constructible Cubic Error Cascade for Approximate Angle Trisection}, preprint, 2026.
\end{thebibliography}
\end{document}
